\begin{answers}{ch:nn}{ニューラルネットワークと誤差逆伝播法}
\item $\calL=6.25$、$\pd{\calL}{w}=-5$、$\pd{\calL}{b}=-5$ です。
グラフは\figref{fig:nn-graph}と同じ形で、$u=wx$、$\hat y=u+b$、$r=\hat y-y$、$\calL=r^2$ の 4 つのノードが並びます。
順伝播の値は $u=-1$、$\hat y=-0.5$、$r=-2.5$、$\calL=6.25$ です。
逆伝播では、まず $\pd{\calL}{r}=2r=-5$ です。足し算と引き算の局所勾配は 1 なので、$\pd{\calL}{\hat y}=\pd{\calL}{u}=\pd{\calL}{b}=-5$ です。
掛け算の局所勾配は相手の値 $x=1$ なので、$\pd{\calL}{w}=1\times(-5)=-5$ です。
直接微分しても $2(wx+b-y)\,x=2\times(-2.5)\times 1=-5$ で、一致します。

\item 次のメソッドを \code{tanh} と同じ形で加えます（\code{import math} は済んでいます）。
\begin{pycode}
def sigmoid(self):
    s = 1.0 / (1.0 + math.exp(-self.data))
    out = Value(s, (self,), "sigmoid")
    def _backward():
        self.grad += s * (1.0 - s) * out.grad  # 局所勾配 × 上流の勾配
    out._backward = _backward
    return out
\end{pycode}
$u=0$ では勾配が $0.25$ になり、中心差分の数値微分（幅 $10^{-6}$）との差は $3\times10^{-11}$ でした。
$\calL=(\sigmoid(wx+b)-1)^2$（$x=1.5$、$w=-0.8$、$b=0.3$）でも、$\pd{\calL}{w}=-0.4383$、$\pd{\calL}{b}=-0.2922$ が、数値微分と $5\times10^{-11}$ 以内で一致しました（実測）。
差が 0 でないのは、数値微分の側の誤差です。

\item 線形モデルと同じになり、訓練損失 0.722・正解率 54\% で頭打ちになります（実測）。
ReLU が無いと、3 枚の全結合層は式\eqref{eq:nn-collapse}のとおり 1 枚の線形の変換にまとまります。
実際、損失は 100 ステップ目までに 0.7217 に下がり、そこから 1,000 ステップまで変わりませんでした（テストでも 52\%）。
これは\secref{sec:nn-linear-limit}の softmax 回帰の結果（損失 0.72、54\%）とほぼ同じです。
ReLU ありの元のスクリプトは、訓練損失 0.0165、正解率 100\% でした。

\item はじめは元より速く下がり、その後、損失が跳ね上がって学習が壊れます（実測）。
損失は 100 ステップ目で 0.041（元のスクリプトは 0.558）まで下がりました。
しかし 120 ステップ目に 5.5 へ跳ね上がり、200 ステップ以降は 1.2〜5.1 の間を行き来して、正解率は 33\% のままでした。
\code{backward()} は勾配を \code{grad} に \code{+=} で足し込みます。
0 に戻さないと、\code{grad} には過去のすべてのステップの勾配の合計が残ります。
最初は勾配の向きがそろっているので、合計が大きくなるほど一歩が大きくなり、速く進みます。
しかし合計はたまる一方なので、やがて一歩が大きすぎて谷を飛び越え、戻れなくなります。

\item 同じにはなりません。重みは乱数の種ごとに違い、隠れ層が小さいと学習に失敗することもあります（実測）。
\file{ch04/mlp_scratch.py} の \code{MLP([2, 2, 2])}（隠れ層の ReLU 2 個、出力は 2 クラスのロジット）を、学習率 0.5 で 1,000 ステップ学習しました。
種 0〜9 のうち正解率 100\% は 1 回だけで、残りは 75\% か 50\%（損失 $\log 2=0.693$ のまま）でした。
ReLU の出力が 4 点すべてで 0 になったニューロンは勾配が 0 になり、学習が止まります。
成功した種 1 では、隠れ層が $h_1=\relu(2.6(x_2-x_1))$、$h_2=\relu(2.6(x_1-x_2))$ になりました。
これは $\abs{x_1-x_2}$ を作る解で、\tabref{tab:nn-xor-hand}の「$h_1-2h_2$」とは別の解き方です。
隠れ層を 4 個にすると 10 回中 6 回、8 個にすると 10 回すべてで 100\% になりました。
XOR を解く重みは無数にあり、学習はそのどれか 1 つにたどり着くだけです。
\end{answers}
